\answer{ex:14:1}{(a)~$0.17$~s e $1.5$~s. (b)~A partir de distâncias menores que $11$~cm.}
\answer{ex:14:2}{O toque sem dor não passa com nenhum $\mu$ enquanto $g\le\beta w_{q\mu}/w_{z\mu}=5$; o toque $\mu=1$ passa com $g>6$: uma sensibilização forte transforma um contato comum em dor.}
\answer{ex:14:3}{(a)~$g^*=(1-k_sC)/(1-k_sA)$, estável com $k_sA<1$. (b)~Sem toque, com $\nu\ge2$; com toque, já com $\nu\ge1.7$.}
\answer{ex:14:4}{$V(t)=-Pe^{-(T-t)/\tau_\gamma}$. (a)~$-Pe^{-T/\tau_\gamma}\approx-0.37P$. (b)~$+Pe^{-(T-t_e)/\tau_\gamma}\approx+0.61P$; ele cresce com $t_e$ até $+P$ e ensina a evitar a dor.}
\answer{ex:14:5}{Com $k_s=4$ não há travamento (o segundo estado estável existe com $k_s\ge4.58$); $T_{\min}\approx4.35$, $1.40$, $0.65$, $0.32\,\tau_s$ com $k_s=4.6$, $5$, $6$, $8$.}
\answer{ex:14:6}{$w_{q\nu}=4/9\approx0.444$, $q_0=2/9\approx0.222$, $w_{q\mu}=49/45\approx1.089$; o toque sem dor é seguro até $g=49/9\approx5.4$.}
\answer{ex:14:7}{(a)~Limiar $\nu_c(g)=\theta/g$ com $\nu_c\ge q_0/w_{q\nu}$, senão $(\theta+\beta q_0)/(g+\beta w_{q\nu})$; $g^*\approx2.45$, $1.0$, $0.25$. (b)~Questão aberta.}
